\subsection{The structure of a conductor equality}
\label{subsec:conductor-equality-structure}

We now prove Theorem~\ref{thm:conductor-equality-structure}. The
positive square just obtained turns a determinant divisor into a
connected curve. Its constant functions split an Euler extension.
This supplies an ambient point through which all the tangent planes
pass, and makes the quotient and its possible singularity explicit.

The exclusion of $a=6$ will be completed in
Proposition~\ref{prop:conductor-six-excluded}; that argument uses
only the structural conclusions and $a\geq5$ proved here.

\begin{proof}[Proof of Theorem~\ref{thm:conductor-equality-structure}]
Put $H=\OO_X(Y)$, $m=pa$ and $d=Y^2$. Let $s_0$ define $Y$.
The equality conclusions of
Propositions~\ref{prop:normalized-mup-conductor}
and~\ref{prop:normalized-mup-singularity-defect} give
\[
 \beta=\omega/s_0^m\in H^0(X,\Omega^1_{X/k}\otimes H^{-m})
\]
whose tangential restriction to $Y$ is nowhere zero. Its zero scheme
$Z$ is finite and disjoint from $Y$: any divisorial zero would meet
the ample curve $Y$. Adjunction gives $K_X\cdot Y=(m-1)d$.
The annihilator of $\beta$ is the saturated invertible subsheaf
\[
 \mathcal F=H^m\otimes\omega_X^{-1}\longrightarrow T_X.
\]
Locally its generator is the Koszul relation between the two
coefficients of $\beta$, which have no common prime divisor.
It is a subbundle outside $Z$, and its normal projection on $Y$
is an isomorphism $\mathcal F|_Y\simeq H|_Y$.

Set
\[
 A=H\otimes\mathcal F^{-1}=\omega_X\otimes H^{-(m-1)},
 \qquad B=A\otimes H.
\]
Thus $A\cdot Y=0$ and $A|_Y\simeq\OO_Y$.
Proposition~\ref{prop:residual-adjoint}, with no additional
divisorial zero, shows that $B$ is nef and big.

\paragraph{Splitting the Euler extension.}
Let $V=H^0(X,H)$ and embed $X$ in $\PP(V^*)$. Pull back the
restricted Euler sequence along
$\mathcal F\to T_{\PP(V^*)}|_X$. This gives a rank-two bundle
$P$ and a commutative diagram
\[
\begin{array}{ccccccccc}
0&\longrightarrow&\OO_X&\longrightarrow&P&\longrightarrow&\mathcal F&\longrightarrow&0\\
&&\Vert&&\downarrow&&\downarrow\\
0&\longrightarrow&\OO_X&\longrightarrow&V^*\otimes H&\longrightarrow&T_{\PP(V^*)}|_X&\longrightarrow&0.
\end{array}
\]
The bundle $P$ remains locally free at $Z$, since it is an extension
of line bundles. Put $E=P^*\otimes H$. There are maps
\[
 0\longrightarrow A\longrightarrow E\longrightarrow H
 \longrightarrow0,\qquad V\otimes\OO_X\longrightarrow E,
\]
and their composite is evaluation. For $s\in V$, denote its
section of $E$ by $e_s$. Locally its components record the value
of $s$ and differentiation in the $\mathcal F$ direction. At a
zero of $s$, the derivative component is independent of the frame.
The section $e_0=e_{s_0}$ is nowhere zero: off $Y$ its value
component is nonzero, and on $Y$ its normal derivative is nonzero.
Consequently
\begin{equation}\label{eq:conductor-euler-extension}
 0\longrightarrow\OO_X\xrightarrow{e_0}E
 \xrightarrow{e_0\wedge-}B\longrightarrow0
\end{equation}
is exact. Its determinant sections $t_s=e_0\wedge e_s$ restrict
on $Y$ to the hyperplane sections $s|_Y$, under the identification
$B|_Y\simeq H|_Y$ given by $e_0$.

Choose $s$ so that $s|_Y$ has $d$ distinct simple zeros. Such
hyperplanes form a nonempty open set, since $Y$ is a smooth
embedded curve. The divisor $C=\operatorname{div}(t_s)$ is
nonempty and does not contain $Y$. Every component meets $Y$;
simple intersection with $Y$ therefore forces every component to
have multiplicity one. A Cartier divisor on a smooth surface has
no embedded components, so $C$ is reduced.

It is also connected. If $C=C_1+C_2$ were a disjoint decomposition
into nonzero effective divisors, then
$C_i^2=B\cdot C_i\geq0$. A zero square would make $C_i$
orthogonal to the positive-square class $B$, and Hodge index
would make $C_i$ numerically zero, contrary to $Y\cdot C_i>0$.
Two positive-square classes cannot be orthogonal either.
Thus $H^0(C,\OO_C)=k$.

On $C$, the image of $e_s$ in the quotient in
\eqref{eq:conductor-euler-extension} vanishes. Hence
$e_s|_C=\lambda e_0|_C$ for a scalar $\lambda\in k$.
Dividing $e_s-\lambda e_0$ by the equation $t_s$ gives a map
$B\to E$ whose composite with $E\to B$ is the identity.
Thus $E\simeq\OO_X\oplus B$.
The surjection $E\to H$ has a nonzero component $B\to H$;
otherwise its image would be generated by $s_0$ and would vanish
on $Y$. This gives a nonzero section of $A^{-1}$.
Its effective zero divisor has intersection zero with the ample
curve $Y$, so it is empty. Therefore
\begin{equation}\label{eq:conductor-linear-canonical}
 A\simeq\OO_X,\qquad B\simeq H,\qquad
 \mathcal F\simeq H,\qquad K_X\sim(m-1)Y.
\end{equation}
This proves linear equivalence, with no injectivity assumption on
restriction of line bundles to $Y$.

\paragraph{The constant tangent point.}
In $E\simeq\OO_X\oplus H$, the component $H\to H$ of the
surjection $E\to H$ is a nonzero scalar. It splits that surjection.
Dualizing and twisting therefore splits the original Euler pullback
compatibly. The inclusion $H\simeq\mathcal F\to T_{\PP(V^*)}|_X$
lifts to a map $H\to V^*\otimes H$, given by a constant nonzero
vector $v\in V^*$. Write $v$ also for its projective point.
The induced tangent direction at $x\ne v$ is that of the line
$xv$. In particular every embedded tangent plane of $X$ contains
$v$. This is the usual notion of a strange embedding
\cite[\S2.1 and Proposition~2.5, author pp.~2--5]{FurukawaStrange}.

The point $v$ lies outside the hyperplane defining $Y$, since the
direction on $Y$ is transverse to that hyperplane. Choose
coordinates with $v=[1:0:\cdots:0]$ and $Y=(x_0=0)\cap X$.
The corresponding section of $T_X\otimes H^{-1}$ has no zero
except possibly $v$. If $v\in X$, its zero ideal is exactly the
maximal ideal of $v$: the ambient coefficients $x_i/x_0$ generate
that ideal, and the tangent inclusion is locally a direct summand.
Thus $Z$ is empty or one reduced point. Let $\epsilon$ be its
length.

At this point choose two ambient affine coordinates $u,w$ forming
regular parameters of $X$. The radial field satisfies
$D(u)=-u$, $D(w)=-w$. The first jet of its annihilating form is
a nonzero scalar multiple of $u\,dw-w\,du$.
Since $v\notin Y$, the same is true of $\omega$.
The derivative of a local unit contributes zero at a zero of the
form. Thus $d\omega(v)=2c\,du\wedge dw$ with $c\ne0$.
Closedness proves
\begin{equation}\label{eq:conductor-epsilon-parity}
 \epsilon=1\quad\Longrightarrow\quad p=2.
\end{equation}
The zero-length formula gives, in both cases,
\begin{equation}\label{eq:conductor-c2X}
 \epsilon=c_2(\Omega_X^1\otimes H^{-m})=c_2(X)+md,
 \qquad c_2(X)=\epsilon-pa\,d.
\end{equation}

\paragraph{The quotient and the normalized torsor.}
Multiplying the twisted field by $x_0$ gives the restriction of
$\delta=x_0\partial/\partial x_0$, with $\delta^p=\delta$.
The homogeneous ideal of $X$ is stable under this derivation.
The eigenspace projections are polynomials in $\delta$, since
$T^p-T$ has distinct roots, so the ideal is graded by weights
modulo $p$. Hence the ambient $\mu_p$ action restricts to $X$.
Its quotient $q:X\to S$ is finite radicial of degree $p$ and
normal. On invariant affine rings, normality follows by
intersection with the invariant function field. The quotient is
projective; equivalently, relative Frobenius factors through it
and makes $S$ finite over $X^{(p)}$.

Let $K=k(X)$ and write the original torsor generically as
$t^p=u$, with $u\notin K^p$. Then $k(N)=L=K(u^{1/p})$.
The field $\delta$ annihilates $du/u$, and hence $u$.
Extending $u$ to a $p$-basis of $K$ shows that
\[
 K^\delta=K^p(u)=L^p.
\]
This identifies the schemes as well as their function fields.
For an affine open $\Spec A\subset S$, let its inverse image in
$N$ be $\Spec C$. Both rings are normal, $C$ is finite over $A$,
and $\operatorname{Frac}A=(\operatorname{Frac}C)^p$.
For $c\in C$, the element $c^p$ is integral over $A$ and lies
in its fraction field, so it belongs to $A$.
Conversely, every $a\in A$ has a $p$th root in
$\operatorname{Frac}C$, integral over $C$, and therefore in $C$.
Thus $A=C^p$. Since $k$ is perfect, these identifications give
\begin{equation}\label{eq:conductor-frobenius-quotient}
 S\simeq N^{(p)},\qquad q\pi=F_{N/k}.
\end{equation}

Away from $v$, the saturated foliation is regular, so its quotient
is smooth. The local result used here is Seshadri's regular
quotient theorem \cite[Proposition~6 and Theorem~2,
pp.~6-11--6-16]{SeshadriCartier}. In rank one, a normalized
generator has $\partial(t)=1$ and $\partial^p=0$; invariant
coordinates give a quotient with regular parameters $t^p,z$.
At $Y$ one can also use the multiplicative coordinates
$\delta(t)=t$, $\delta(z)=0$.

If $v\in X$, then $p=2$ and its completed invariant ring is
\[
 k[[u^2,uw,w^2]]\simeq k[[a,b,c]]/(ac-b^2).
\]
Indeed, invariant monomials have even total degree, and are
products of $u^2,uw,w^2$, with the displayed sole relation.
Invariants commute with completion through the weight-zero
direct summand. This is an $A_1$ singularity. Blowing up $v$
in $X$, with exceptional curve $R$, resolves the quotient:
on the chart $w=uz$ the lifted field satisfies
$\delta(u)=-u$, $\delta(z)=0$, and the quotient chart has
coordinates $u^2,z$. The other chart is analogous.
Let $q':X'\to T'$ be this smooth quotient and $\rho:T'\to S$
the induced resolution. Its exceptional curve $E$ satisfies
$q'^*E=2R$ and $E^2=-2$.
It is the minimal resolution. It is crepant as well: since $S$
is Gorenstein, write $K_{T'}=\rho^*K_S+cE$; adjunction gives
$K_{T'}\cdot E=0$, whence $c=0$.
For $\epsilon=0$ set $X'=X$ and $T'=S$.
Equation~\eqref{eq:conductor-frobenius-quotient} transfers these
smoothness and singularity statements to $N$; its minimal
resolution $T$ has Frobenius twist $T'$.

\paragraph{Canonical class and Chern numbers.}
The normalized conductor formula gives
\[
 K_N=\pi^*(K_X-a(p-1)Y)\sim(a-1)\pi^*Y.
\]
For $a\geq2$ it is ample. The crepant resolution has nef and
big canonical divisor, so $T$ and $T'$ are minimal surfaces of
general type. Their canonical squares are
\[
 K_T^2=K_{T'}^2=K_N^2=p(a-1)^2d.
\]

We compute their second Chern numbers directly on smooth surfaces.
For a regular rank-one foliation $F\subset T_{X'}$, let $Q$
be its annihilator in $\Omega^1_{X'}$.
The invariant coordinates just used give exact sequences
\[
 0\to Q\to\Omega^1_{X'}\to F^{-1}\to0,
 \qquad
 0\to F^{-p}\to q'^*\Omega^1_{T'}\to Q\to0.
\]
Here $Q=\omega_{X'}\otimes F$; the $p$th power in the second
kernel follows from the transition of the local coordinate $t^p$.
Writing $f=c_1(F)$, determinants and second Chern classes give
\[
 q'^*K_{T'}=K_{X'}-(p-1)f,\qquad
 c_2(X')=-f\cdot(K_{X'}+f),\qquad c_2(T')=c_2(X').
\]
For $\epsilon=0$, $f=Y$. For $\epsilon=1$, the lifted
field has fixed divisor $b^*Y+R$, so $f=b^*Y+R$ and
$K_{X'}=b^*K_X+R$.
Using $R^2=-1$ and $b^*Y\cdot R=0$, one obtains
\begin{equation}\label{eq:conductor-structure-invariants}
\begin{gathered}
 K_T^2=p(a-1)^2d,\qquad c_2(T)=-pa\,d+2\epsilon,\\
 12\chi(\OO_T)=p(a^2-3a+1)d+2\epsilon.
\end{gathered}
\end{equation}
The last equality is Noether's formula.
As a separate check, Noether's formula on $X$ and
Corollary~\ref{cor:normalized-equality-euler} give
$12\chi(\OO_N)=p(a^2-3a+1)d+p\epsilon$;
equation~\eqref{eq:conductor-epsilon-parity} makes the two
expressions equal.

\paragraph{The first four conductors are impossible.}
If $a=1$, $K_T$ is numerically trivial.
The classification of smooth projective surfaces with numerically
trivial canonical divisor gives $c_2(T)\geq0$
\cite[\S2.1, Table~1, author p.~3]{MatsumotoMuN2023}.
The table has Betti pairs $(b_1,b_2)=(4,6),(0,22),(0,10),(2,2)$,
giving Euler numbers $0,24,12,0$.
For $\epsilon=0$, this contradicts $c_2(T)=-pd<0$.
For $\epsilon=1$, $p=2$ and $2-2d\geq0$, hence $d=1$.
But $H|_Y$ is very ample of degree one, so $Y$ is a line.
This contradicts $2g(Y)-2=pa\,d>0$.

For $a\geq2$, Gu--Sun--Zhou's inequality applies to the smooth
minimal surface $T$ of general type:
\[
 K_T^2\leq32\chi(\OO_T)
 \qquad\text{\cite[Theorem~1.5, p.~614;
 Corollary~3.4, pp.~623--624]{GuSunZhou}.}
\]
In particular $\chi(\OO_T)>0$.
If $a=2$, equation~\eqref{eq:conductor-structure-invariants}
gives $12\chi(\OO_T)=-pd+2\epsilon\leq0$, a contradiction.
If $a=3$ and $\epsilon=0$, it gives
$K_T^2=4pd>32pd/12=32\chi(\OO_T)$.
If $a=3$ and $\epsilon=1$, it gives
\[
 K_T^2=8d,\qquad \chi(\OO_T)=(d+1)/6.
\]
The inequality forces $d\leq2$, but neither positive integer
$d=1,2$ makes this Euler characteristic integral.
Thus $a\geq4$, proving ampleness and minimality.

We have established $a\geq4$. Suppose $a=4$, and let
$\rho:T\to N$ be the minimal resolution, with $\rho$ the identity
when $N$ is smooth. Its canonical identity and crepancy give
\[
 K_T\sim3A,\qquad A=\rho^*\pi^*Y,
\]
where $A$ is an integral Cartier divisor class. The invariant
formulas give
\[
 K_T^2=9pd,\qquad 12\chi(\OO_T)=5pd+2\epsilon.
\]
Since $\epsilon=1$ implies $p=2$, it follows that
\[
 12\bigl(K_T^2-18\chi(\OO_T)\bigr)
 =18pd-36\epsilon=18p(d-\epsilon)>0.
\]
For $\epsilon=0$, positivity follows from $d>0$.
For $\epsilon=1$, the value $d=1$ would make the very ample
curve $Y$ a line, contradicting $2g(Y)-2=pa\,d>0$.

The surface $T$ is smooth minimal of general type.
Gu--Sun--Zhou's theorem therefore gives a genus-two Albanese
fibration
\cite[Theorem~1.5, p.~614; Corollary~3.4 and its proof,
pp.~623--624]{GuSunZhou}.
For a fibre divisor $F$, one has $F^2=0$ and
$K_T\cdot F=2$ by adjunction.
The genus is the arithmetic genus in the flat-fibration
convention of that theorem, so this calculation does not require
a smooth fibre. It contradicts $K_T\sim3A$, since
$A\cdot F$ is an integer.
Therefore $a\geq5$.

\paragraph{The divisibility of the degree.}
If $v\notin X$, projection from $v$ is a morphism defined by
$x_1,\ldots,x_r$, with pullback of $\OO(1)$ equal to $H$.
Its coordinate ratios are invariant under $\mu_p$, so it factors
through $q:X\to S$. For the descended line bundle $L$ on $S$,
$q^*L=H$. Hence $d=pL^2$.
If $v\in X$, the base ideal of the projection is the maximal
ideal of that reduced point. On $X'$ the resolved projection
has line bundle $b^*H-R$ and factors through $q':X'\to T'$.
Its descended line bundle satisfies $q'^*L=b^*H-R$; thus
$d-1=pL^2$. This proves $p\mid(d-\epsilon)$.
No very ampleness of the descended line bundle is required.
\end{proof}

\begin{corollary}\label{cor:conductor-not-complete-intersection}
In Theorem~\ref{thm:conductor-equality-structure}, the embedding
of $X$ by $\OO_X(Y)$ is not a complete intersection.
\end{corollary}

\begin{proof}
The embedding is strange, as the proof shows, and is not linear
since $K_X\cdot Y=(pa-1)d>0$.
Furukawa's theorem states that a nonlinear smooth complete
intersection is strange exactly when it is an odd-dimensional
quadric in characteristic two
\cite[Theorem~1.1, author p.~2, proved in \S3]{FurukawaStrange}.
Our surface has even dimension, so it cannot be such a complete
intersection.
\end{proof}
